\section{The Phonon Sector (Neutrinos)}

The blind spectral scan in Section V searched for linear harmonics ($k \cdot \Delta^n$), successfully identifying the ``Knots'' of the lattice. It returned a null result for the neutrino sector because the generator was not configured to search the inverse-harmonic domain ($1/\Delta^n$). This distinction is not a failure of the model, but a confirmation of the \textbf{Structural Duality} of the substrate.

While charged fermions are \textbf{Topological Solitons} (``Knots'') defined by the dimensional drag ($\pi D$), the neutrino sector represents a fundamentally distinct class of geometric excitation: \textbf{Lattice Phonons}. Standard fermions exist \textit{in} the lattice; Neutrinos are vibrations \textit{of} the lattice. Consequently, their mass spectrum follows an \textbf{Inverse Resonance Hierarchy} ($M \propto \Delta^{-n}$), and their effective mass is context-dependent.

We posit a \textbf{Context-Dependent Mass Operator}: Neutrinos exhibit a light \textbf{Rest Mass} ($M_{grav}$) when decoupled, and a heavy \textbf{Interaction Mass} ($M_{weak}$) when coupled to the gauge field.

\textbf{Structural Note (Spin-Statistics Compliance):} We use the term ``Phonon'' to denote the inverse-harmonic scaling ($1/\Delta^n$) of a delocalized state. However, unlike acoustic phonons in condensed matter (which are scalar bosons), neutrino excitations propagate through the \textbf{Chiral State Space} ($\nu=16$). Because the substrate nodes are fundamental spinors (as derived in Paper I/V), their collective vibrations are \textbf{Spinor Waves}. Consequently, they preserve Fermionic statistics ($s=1/2$) and strictly satisfy the Pauli Exclusion Principle.

\subsection{The Gravitational Mass (The Rest State)}

The Standard Model faces a tension where oscillation experiments imply $\Sigma m_\nu \gtrsim 60$ meV, while cosmological surveys constrain $\Sigma m_\nu < 120$ meV. We resolve this by identifying neutrinos as lattice phonons.

\subsubsection{The Gravitational Anchor (\texorpdfstring{$M_{rest}$}{Mrest})}
In Paper I (Section XIII.C), we derived the \textbf{Lattice Floor ($M_{min}$)}, the minimum resolvable vibration of the substrate defined by the thermal noise limit:
\begin{equation}
    M_{min} = m_e \cdot \frac{\pi \alpha}{\nu \Delta^3} \approx 9.2 \text{ meV}
\end{equation}
We identify this geometric ground state physically as the \textbf{Heaviest Neutrino Eigenstate ($\nu_3$)}.

\subsubsection{The Linear Phonon Exclusion}
The inverse resonance series theoretically begins at $n=1$. However, a linear phonon corresponds to a mass scale of $M_{n=1} = M_{min} \cdot \Delta \approx 0.4 \text{ eV}$. 

This ``Linear Phonon'' represents a sterile fluctuation. Unlike areal ($n=2$) and volumetric ($n=3$) phonons which can couple to the 2D boundary and 3D bulk of the lattice topology, a 1D phonon lacks the spatial dimensionality to couple effectively to the gauge fields. Consequently, the $n=1$ state is strictly sterile (decoupled) or topologically unstable, forcing the observable spectrum to begin at the Areal Tier.

\subsubsection{The Inverse Resonance Hierarchy}
Because neutrinos are inverse resonances (phonons), they occupy the harmonic tiers $1/\Delta^n$. This creates a forced \textbf{Normal Hierarchy} relative to the base lattice floor:

\begin{itemize}
    \item $\nu_3$ (Volume Mode / Base): $M_{min} \times \Delta^0 \approx \mathbf{9.2 \text{ meV}}$
    \item $\nu_2$ (Area Mode): $M_{min} \times \Delta^{-1} \approx \mathbf{0.21 \text{ meV}}$
    \item $\nu_1$ (Linear Mode): $M_{min} \times \Delta^{-2} \approx \mathbf{0.005 \text{ meV}}$
\end{itemize}

\textbf{Total Rest Mass:}
\begin{equation}
    \Sigma m_\nu = 9.2 + 0.21 + 0.005 \approx \mathbf{9.4 \text{ meV}}
\end{equation}

\textbf{Conclusion:} Cosmology measures this Rest Mass. This explains why Planck and DES data prefer $\Sigma m_\nu \approx 0$ (effectively indistinguishable from zero within current error bars), while oscillation experiments measure the heavier Interaction Mass (See Section B).

\subsection{The Interaction Mass (The Weak State)}

Why do oscillation experiments measure mass splittings corresponding to $\sim 50$ meV if the rest mass is $\sim 9$ meV?

\textbf{1. Derivation of the Refractive Enhancement Factor ($\eta_W$)} \\
The Interaction Mass enhancement arises from the neutrino's coupling to the $SU(2)_L$ gauge boson. Because neutrinos are phonons, they couple through the \textbf{Lattice Density of States}, which is enhanced by the \textbf{Interaction Aperture}:
\begin{equation}
    \eta_W = \sigma + 1 = 5 + 1 = \mathbf{6}
\end{equation}

\textbf{Validation:} This aperture factor is not arbitrary; it is the same geometric constraint that governs the $Z$-boson decay width. As demonstrated in \textbf{Paper III, Section V.E}, the ratio of $W$-to-$Z$ decay widths is determined by the aperture ratio $5/6$, matching experiment to $99.8\%$. We therefore apply this validated factor $\eta_W=6$ to the neutrino refractive index.

\textbf{2. The Refractive Index Mechanism} \\
When a neutrino participates in a Weak Interaction (flavor oscillation or beta decay), it propagates through the dense ``Protocol Layer'' of the gauge field. Just as a photon acquires an effective mass in a plasma, the neutrino phonon acquires an \textbf{Interaction Mass} due to Protocol Latency, scaled by $\eta_W$.
\begin{equation}
    M_{weak} \approx M_{rest} \times \eta_W \approx 9.2 \text{ meV} \times 6 \approx \mathbf{55 \text{ meV}}
\end{equation}

\textbf{3. Consistency with Oscillations}
\begin{itemize}
    \item \textbf{Atmospheric Splitting:} $\Delta m^2_{23} \approx (55 \text{ meV})^2 - (1.2 \text{ meV})^2 \approx 3.0 \times 10^{-3} \text{ eV}^2$. (Consistent with observed $\approx 2.5 \times 10^{-3} \text{ eV}^2$).
    \item \textbf{Open Problem:} This simple scaling predicts a solar splitting $\Delta m^2_{12}$ smaller than observed, suggesting the refractive index is generation-dependent. However, the \textit{scale} of the atmospheric mass is correctly reproduced without free parameters.
\end{itemize}

\subsection{The Decoupling Limit (Thermal History)}

The transition between measuring $M_{weak}$ and $M_{rest}$ is governed by the Interaction State. The neutrino decoupling temperature is defined geometrically as:
\begin{equation}
    T_{decouple} \approx (G_F^2 M_{Planck})^{-1/3} \approx 1 \text{ MeV}
\end{equation}

\begin{itemize}
    \item \textbf{Coupled Regime ($T > 1$ MeV):} In the early universe or inside reactors/stars, $\Gamma_{weak} > H(t)$. Neutrinos are tightly coupled to the gauge field and exhibit the heavy Interaction Mass ($M_{weak} \approx 55$ meV).
    \item \textbf{Decoupled Regime ($T < 1$ MeV):} In the expanding universe, once the temperature drops below 1 MeV, the weak interaction rate falls below the expansion rate. Neutrinos ``freeze out'' into free lattice phonons, exhibiting only their intrinsic Rest Mass ($M_{rest} \approx 9$ meV).
\end{itemize}